\documentclass[a4paper]{article}

%% Language and font encodings
\usepackage[english]{babel}
\usepackage[utf8]{inputenc}
\usepackage[T1]{fontenc}
\usepackage{amsfonts}

%% Sets page size and margins
\usepackage[a4paper,top=3cm,bottom=2cm,left=3cm,right=3cm,marginparwidth=1.75cm]{geometry}

%% Useful packages
\usepackage{amsmath}
\usepackage{amssymb}
\usepackage{amsthm}
\usepackage{graphicx}
\usepackage[section]{placeins}
\usepackage[colorinlistoftodos]{todonotes}
\usepackage[colorlinks=true, allcolors=blue]{hyperref}

%% Table packages
\usepackage{array}
\newcolumntype{P}[1]{>{\centering\arraybackslash}p{#1}}
\usepackage{multirow}

\title{Mathematical and Physical Modeling of Entangled Photon Propagation in Optical Fibers for MATLAB Simulation}
\author{Marco Cioci\\
Politecnico di Milano\\
Student ID: 10694772\\
\texttt{marco.cioci@mail.polimi.it}}
\date{}

\begin{document}
\maketitle

\tableofcontents

%========================================================
\newpage
\section{General Framework}

\subsection{Purpose and Scope}
The purpose of this document is to define a coherent working framework for the mathematical modeling and numerical simulation of polarization-entangled photon pairs propagating through optical fibers.

This text is not intended to appear directly in the final thesis manuscript. Instead, it serves as an internal technical guide throughout the development of the project. Its role is to support communication, maintain consistency of notation and assumptions, and provide a stable reference for the progressive validation of both the theoretical model and the MATLAB implementation.

The central objective is to organize the physical assumptions, mathematical formalism, and simulation strategy into a reproducible structure, so that each computational step can be traced back to a precise theoretical element. In this way, the simulator is not developed as an isolated numerical tool, but as a direct operational translation of the underlying quantum model.

\paragraph{Simulation notes.}
At the MATLAB level, this document should be treated as the reference specification for notation, model hierarchy, and implementation priorities. The simulator should therefore be organized into clearly identifiable modules, each corresponding to a specific theoretical block. This modular structure is essential both for step-by-step validation of intermediate results and for future extensions of the model.

\subsection{Theoretical Formalism and Notation}

The system is described within the framework of quantum information theory. The relevant mathematical structures include finite-dimensional Hilbert spaces, quantum states, tensor products, unitary operators, density operators, and projective measurements. Elements of classical optics are introduced only when required to justify the physical origin of the effective quantum transformations.

The notation and formal tools are aligned, as far as possible, with the framework of \textit{Quantum Computation and Quantum Information} by Nielsen and Chuang.

In this formalism:
\begin{itemize}
    \item single subsystems are represented by vectors in \(\mathbb{C}^2\),
    \item composite systems are described by tensor-product spaces,
    \item evolutions are represented by linear operators (unitary in the ideal case),
    \item measurements are described through Hermitian operators.
\end{itemize}

All states and operators are represented as vectors and matrices over \(\mathbb{C}\), with dimensions determined by the tensor-product structure of the system.

\paragraph{Simulation notes.}
The simulator adopts a quantum-information-oriented representation from the outset. Pure states are represented as complex vectors, mixed states as density matrices, and evolutions and measurements as matrix operators.

This abstract structure defines the numerical backbone of the simulator and is independent of the specific physical encoding used.


\subsection{Physical Scenario and Experimental Setting}

The physical setting considered in this work is the following. A spontaneous parametric down-conversion (SPDC) source generates pairs of polarization-entangled photons, ideally prepared in a Bell state. The two photons are injected into distinct optical paths, denoted \(A\) and \(B\), and propagate through independent optical fibers.

Each fiber acts on the polarization degree of freedom of the corresponding photon. In a general description, this action can be represented by a unitary operator acting on a single-qubit Hilbert space. However, motivated by experimental constraints and by the limited degree of control available on the full polarization transformation, the model is initially restricted to a simplified effective description.

After compensation of global polarization rotations, the dominant residual effect can be modeled as a relative phase shift between the horizontal and vertical polarization components. In the qubit representation, this corresponds to a rotation around the computational \(Z\)-axis of the Bloch sphere.

The objective at this stage is therefore not to reproduce the full electromagnetic propagation in optical fibers, but to construct an effective quantum model that captures the relevant impact of fiber propagation on entanglement and measurable correlations.

\paragraph{Simulation notes.}
At the simulator level, the physical scenario is reduced to the minimal set of elements required for computation: an initial entangled two-qubit state, two labeled channels \(A\) and \(B\), and one local transformation per channel.

In the first implementation, unnecessary optical detail is neglected and each fiber is described exclusively through its effective action on polarization. This corresponds to a local unitary model of the form
\[
U_A \otimes U_B.
\]

In the currently implemented baseline experiment, the evolution is further simplified to
\[
U_A = U(\theta),
\qquad
U_B = I,
\]
so that the full state evolution depends on a single phase parameter \(\theta\). This configuration defines the first validation scenario of the simulator, allowing direct comparison between analytical predictions and numerical results for correlation observables.

%========================================================
\section{Modeling Strategy}

\subsection{Progressive Modeling Approach}

The modeling approach follows a progressive refinement strategy, in which the system is described through successive levels of approximation of increasing physical complexity.

Each stage introduces additional physical effects while remaining consistent with the previous ones. This ensures analytical transparency and preserves continuity between simplified and refined models.

From a theoretical perspective, this approach isolates the contribution of individual mechanisms. From a computational perspective, it enables a layered simulator structure, where each refinement corresponds to a well-defined extension of the existing codebase.

%--------------------------------------------------------

\subsection{Hierarchy of Modeling Levels}

The modeling workflow is structured as a sequence of progressively refined descriptions:
\[
\begin{array}{c}
\text{Ideal Bell State} \\
\downarrow \\
\text{Local Phase Model (Z-rotations)} \\
\downarrow \\
\text{General Local Unitary Transformations} \\
\downarrow \\
\text{Mixed-State Description (Density Operators)} \\
\downarrow \\
\text{Noisy Quantum Channels} \\
\downarrow \\
\text{Effective Fiber Models}
\end{array}
\]

More explicitly:

\begin{itemize}
    \item \textbf{Level 1: Minimal phase model} \\
    Pure, maximally entangled state. Fiber action reduced to a relative phase. Analytical dependence of correlations on a single parameter.

    \item \textbf{Level 2: General unitary dynamics} \\
    Arbitrary polarization rotations described by local operators \(U_A \otimes U_B\).

    \item \textbf{Level 3: Mixed-state description} \\
    Density operators used to represent statistical mixtures and loss of coherence.

    \item \textbf{Level 4: Noisy quantum channels} \\
    Effective completely positive trace-preserving (CPTP) maps modeling depolarization and decoherence.

    \item \textbf{Level 5: Effective fiber models} \\
    Physically motivated descriptions, such as concatenations of random unitary segments or stochastic polarization evolution.
\end{itemize}

This hierarchy provides both a conceptual and an implementation roadmap. The initial levels remain analytically tractable and allow direct validation, while the later levels increase physical realism and connect to experimentally relevant effects.

\paragraph{Simulation notes.}
This hierarchy should be mirrored explicitly in the codebase. Each level should correspond to a distinct layer of functionality: state preparation, local unitary evolution, density-matrix handling, channel application, and advanced propagation models. This structure ensures that each refinement appears as an additional computational step rather than a structural modification of the simulator.

%--------------------------------------------------------

\subsection{Baseline Model: Pure-State Description}

At the first stage, the bipartite system is described as a pure state in \(\mathbb{C}^2 \otimes \mathbb{C}^2\), and fiber propagation is modeled through local unitary transformations. Particular emphasis is placed on relative phase shifts induced by the fibers.

The main objective is to analyze how these phase shifts affect measurable two-photon correlations. The relevant observables are operators of the form
\[
\sigma_i \otimes \sigma_j,
\qquad i,j \in \{x,y,z\},
\]
as well as operators corresponding to arbitrary measurement directions on the Bloch sphere.

This baseline model is the natural entry point because it allows explicit analytical expressions, direct visualization of state evolution, and immediate comparison between analytical predictions and numerical results.

\paragraph{Simulation notes.}
At this level, the simulator operates entirely on pure two-qubit states represented as \(4 \times 1\) complex vectors. The core tasks are: generation of the Bell state, application of local operators \(U_A\) and \(U_B\), and evaluation of expectation values.

In the current implementation, the standard aligned correlations
\[
\langle \sigma_x \otimes \sigma_x \rangle,\qquad
\langle \sigma_y \otimes \sigma_y \rangle,\qquad
\langle \sigma_z \otimes \sigma_z \rangle
\]
are computed and analyzed as functions of the phase parameter \(\theta\). This constitutes the first validation stage of the simulator.

%--------------------------------------------------------
\subsection{Extension to Mixed-State Formalism}

At the next level, the model incorporates statistical effects and decoherence. In this regime, the pure-state description is no longer sufficient and must be replaced by a density-operator formalism.

This framework enables the description of:
\begin{itemize}
    \item reduced states of individual photons,
    \item loss of coherence due to unresolved degrees of freedom,
    \item effective non-unitary dynamics arising from environmental interactions or statistical averaging.
\end{itemize}

It also allows the computation of physically relevant quantities such as fidelity, concurrence, and other entanglement measures.

The introduction of density operators is therefore required whenever the system is described as an ensemble or when partial information about the system is available.

An important derived quantity is the reduced density operator of a subsystem, obtained via partial trace. For instance, the reduced state of subsystem \(A\) is
\[
\rho_A = \mathrm{Tr}_B(\rho_{AB}).
\]

For maximally entangled states such as \( |\Psi^+\rangle \), the reduced density operators are maximally mixed:
\[
\rho_A = \rho_B = \frac{1}{2} I.
\]

This property reflects the fact that each individual photon carries no definite polarization, even though the joint state exhibits strong correlations. It provides a direct connection between entanglement and the absence of local information, and plays a central role in the interpretation of both deterministic and ensemble-based simulations.

Expectation values in the density-operator formalism are computed as
\[
\langle O \rangle = \mathrm{Tr}(\rho O),
\]
which generalizes the pure-state expression \( \langle \psi | O | \psi \rangle \). This formulation is essential when describing ensemble-averaged systems.

\paragraph{Simulation notes.}
The simulator must be extended so that all relevant routines operate on density matrices. This includes \(4 \times 4\) matrix representations, partial trace operations, quantum channel application, and metric evaluation.

The pure-state implementation should be preserved as a special case of this more general framework, ensuring consistency between deterministic and ensemble-based simulations.

%--------------------------------------------------------

\subsection{Simulator Design and Architecture}

The simulator is designed as a modular implementation of the theoretical model.

The initial modules include:
\begin{itemize}
    \item Bell-state generation,
    \item local unitary propagation,
    \item correlation evaluation.
\end{itemize}

More advanced modules introduce:
\begin{itemize}
    \item quantum channels,
    \item reduced density operators,
    \item entanglement measures,
    \item experimentally relevant observables such as Bell parameters and Stokes-like quantities.
\end{itemize}

The theoretical structure is directly reflected in the computational architecture, ensuring consistency, traceability, and extensibility.

The implementation distinguishes clearly between reusable computational kernels and higher-level experiment scripts, allowing the simulator to evolve into a structured and reusable research tool.

\paragraph{Simulation notes.}
At the current stage, this modular design is reflected in the directory structure:
\begin{itemize}
    \item \texttt{Core}: state vectors, operators, expectation values, correlation routines;
    \item \texttt{Models}: physical propagation models;
    \item \texttt{Utils}: validation routines and auxiliary tools;
    \item \texttt{Experiments}: scripts implementing complete simulation studies.
\end{itemize}

This separation enforces a strict distinction between general mathematical operations and physically specific modeling choices.

%========================================================
\section{Physical and Mathematical Representation}

\subsection{Quantum System and Hilbert Space Structure}

The physical system considered consists of polarization-entangled photon pairs generated by spontaneous parametric down-conversion (SPDC). The two photons propagate along distinct paths, denoted \(A\) and \(B\), forming a bipartite system.

In this work, the qubit is encoded in the polarization degree of freedom:
\[
|H\rangle \equiv |0\rangle,
\qquad
|V\rangle \equiv |1\rangle.
\]

Accordingly, the computational basis of the bipartite system is
\[
\{
|00\rangle,\,
|01\rangle,\,
|10\rangle,\,
|11\rangle
\},
\]
corresponding to
\[
\{
|H\rangle_A |H\rangle_B,\,
|H\rangle_A |V\rangle_B,\,
|V\rangle_A |H\rangle_B,\,
|V\rangle_A |V\rangle_B
\}.
\]

The Hilbert space of the total system is therefore
\[
\mathcal{H}_{AB} = \mathcal{H}_A \otimes \mathcal{H}_B
\cong \mathbb{C}^2 \otimes \mathbb{C}^2,
\]
which has dimension \(4\).

As a consequence:
\begin{itemize}
    \item pure states are represented by complex vectors of length \(4\),
    \item operators acting on the full system are \(4 \times 4\) matrices.
\end{itemize}

This construction specifies the concrete realization of the abstract quantum-information framework introduced previously.

\paragraph{Simulation notes.}
At MATLAB level, bipartite pure states are represented as \(4 \times 1\) complex vectors and operators as \(4 \times 4\) matrices. The basis ordering is fixed as
\[
\{|00\rangle,|01\rangle,|10\rangle,|11\rangle\},
\]
and must remain invariant, since all tensor products, observables, and density matrices depend on this convention.

This ordering is already enforced in the current codebase and must be preserved in all future extensions.

%--------------------------------------------------------

\subsection{Bell-State Initialization}

As a reference initial condition, we consider the Bell state
\[
|\Psi^{+}\rangle
=
\frac{1}{\sqrt{2}}
\left(
|H\rangle_A |V\rangle_B
+
|V\rangle_A |H\rangle_B
\right),
\]
which, in computational notation, is
\[
|\Psi^{+}\rangle
=
\frac{1}{\sqrt{2}}
\left(
|01\rangle + |10\rangle
\right).
\]

In the ordered computational basis
\[
\{|00\rangle, |01\rangle, |10\rangle, |11\rangle\},
\]
this state is represented by the column vector
\[
|\Psi^{+}\rangle
=
\frac{1}{\sqrt{2}}
\begin{pmatrix}
0\\
1\\
1\\
0
\end{pmatrix}.
\]

This state is maximally entangled and provides the natural input for the baseline version of the simulator. It also serves as the reference state with respect to which later quantities such as fidelity and Bell-inequality violation may be evaluated.

The choice of a Bell state is especially convenient because it provides both a clean theoretical benchmark and a physically relevant entangled resource. It also allows the effect of fiber propagation to be interpreted directly in terms of coherence and two-photon correlation changes.

\paragraph{Simulation notes.}
The simulator should include a dedicated state-preparation routine returning the Bell-state vector in canonical basis ordering. This reference state is useful for normalization checks, debugging, comparison with analytical formulas, and later evaluation of propagated-state fidelity.

At the current stage, the simulator already includes a Bell-state generator returning the four standard Bell states:
\[
|\Psi^{\pm}\rangle
=
\frac{1}{\sqrt{2}}\left(|01\rangle \pm |10\rangle\right),
\qquad
|\Phi^{\pm}\rangle
=
\frac{1}{\sqrt{2}}\left(|00\rangle \pm |11\rangle\right),
\]
although the present baseline model uses \( |\Psi^+\rangle \) as the default reference input.

%--------------------------------------------------------

\subsection{Fiber Propagation as Local Unitary Evolution}

Each photon propagates through its own optical fiber. At the level of the polarization qubit, the action of each fiber is described by a linear operator acting locally on the corresponding single-photon Hilbert space. In the most general ideal lossless description, this operator is unitary. Therefore, if \(U_A\) and \(U_B\) denote the single-qubit operators associated with the two fibers, the evolution of the initial two-photon state is
\[
|\psi_{\mathrm{out}}\rangle
=
(U_A \otimes U_B)\, |\psi_{\mathrm{in}}\rangle.
\]

This is the natural bipartite extension of single-qubit evolution and follows directly from the tensor-product structure of composite quantum systems. It is the appropriate starting point for later refinements, including random unitary perturbations, concatenated fiber sections, and effective channel descriptions.

In this framework, the two fibers act independently on the two polarization subsystems, and the total evolution is obtained by tensor composition of the corresponding local transformations. This description is valid as long as the fibers do not induce any direct interaction between the two photons.

\paragraph{Simulation notes.}
The fundamental evolution primitive in MATLAB should be the tensor-product action
\[
\psi_{\mathrm{out}} = (U_A \otimes U_B)\psi_{\mathrm{in}}.
\]
It is convenient to isolate this in a dedicated routine, so that any later modification of the single-fiber model automatically propagates to the full bipartite simulation.

This abstraction is already present in the implementation, where local single-qubit matrices are combined into the global operator \(U_A \otimes U_B\) before being applied to the input state.

%--------------------------------------------------------

\subsection{Reduced Phase Model}

Although a general fiber transformation may involve arbitrary polarization rotations, birefringence, and mode-coupling effects, the first approximation adopted here is simpler and directly motivated by the experimental control scheme. After partial compensation of polarization drifts, the dominant residual effect is modeled as a relative phase shift between the horizontal and vertical components, corresponding to an effective rotation around the \(Z\)-axis of the Bloch sphere.

At the single-qubit level, this is represented by the diagonal unitary
\[
U(\theta)
=
\begin{pmatrix}
1 & 0\\
0 & e^{i\theta}
\end{pmatrix},
\]
up to an irrelevant global phase. If the two fibers induce phases \(\theta_A\) and \(\theta_B\), then the total transformation is
\[
U_A \otimes U_B
=
U(\theta_A)\otimes U(\theta_B).
\]

Applying this operator to the reference Bell state \( |\Psi^+\rangle \) gives
\[
|\psi(\theta_A,\theta_B)\rangle
=
\frac{1}{\sqrt{2}}
\left(
e^{i\theta_B}|01\rangle
+
e^{i\theta_A}|10\rangle
\right).
\]
By factoring out a global phase, which has no effect on measurement statistics, this state can be rewritten as
\[
|\psi(\theta)\rangle
=
\frac{1}{\sqrt{2}}
\left(
|01\rangle
+
e^{i\theta}|10\rangle
\right),
\]
where
\[
\theta = \theta_A - \theta_B.
\]

Thus, in the reduced model, the physically relevant fiber action is encoded in a single parameter: the relative phase \(\theta\) between the two components of the entangled state. This parameter depends on propagation conditions such as birefringence, path length, and environmental fluctuations, and it becomes the main control variable in the first stage of the simulator.

This reduced description is useful because it isolates the dominant effect of interest while preserving the full two-qubit entangled structure.

\paragraph{Simulation notes.}
This is the first concrete propagation model to implement. The simulator should allow \(\theta_A\) and \(\theta_B\) as input parameters, while also recognizing that the observable behavior depends only on \(\theta=\theta_A-\theta_B\). In practice, both forms are useful: the two-phase form for physical interpretation, and the reduced one-parameter form for analytical checks and parameter sweeps.

At the current implementation stage, the reduced model is realized through a dedicated single-qubit phase-unitary routine and is used operationally in the first complete experiment by fixing one arm to the identity and sweeping the phase in the other arm.

%--------------------------------------------------------

\subsection{Interpretation and Role of the Reduced Model}

The reduced phase model is particularly useful for three reasons. First, it preserves a pure-state description, keeping the initial implementation analytically transparent. Second, it isolates the effect most directly relevant for the first experimental analysis, namely the deformation of two-photon correlations induced by a relative phase. Third, it provides the natural starting point for later extensions to density operators and noisy channels without changing the basic bipartite structure of the problem.

For this reason, the first stage of the simulator is built around the family of states
\[
|\psi(\theta)\rangle
=
\frac{1}{\sqrt{2}}
\left(
|01\rangle + e^{i\theta}|10\rangle
\right),
\]
from which one computes expectation values of local and nonlocal observables, correlation functions, and Bell-type quantities as functions of \(\theta\).

The reduced model therefore plays a dual role: it is both a physically motivated approximation to fiber propagation and the minimal operational model from which one can derive clean and testable predictions.

\paragraph{Simulation notes.}
The first validated MATLAB study should therefore consist of a parameter sweep in \(\theta\), with numerical evaluation of state vectors, expectation values, and correlation quantities. This provides a minimal but complete simulator: one input family of states, one effective propagation parameter, and a set of measurable outputs that can be directly plotted and compared with theory.

At the current stage, this is implemented through a sweep routine that returns a structured collection of sampled phase values and corresponding correlation outputs, together with a plotting utility for visualizing selected correlations as functions of \(\theta\).

%--------------------------------------------------------

\subsection{Correlation Observables: Theory and Implementation}

At the present stage of development, the observables explicitly implemented in the simulator are the standard aligned two-qubit Pauli correlations:
\[
\sigma_x \otimes \sigma_x,\qquad
\sigma_y \otimes \sigma_y,\qquad
\sigma_z \otimes \sigma_z.
\]

These are constructed from the single-qubit Pauli matrices
\[
\sigma_x=
\begin{pmatrix}
0 & 1\\
1 & 0
\end{pmatrix},
\qquad
\sigma_y=
\begin{pmatrix}
0 & -i\\
i & 0
\end{pmatrix},
\qquad
\sigma_z=
\begin{pmatrix}
1 & 0\\
0 & -1
\end{pmatrix}.
\]

Given a pure two-qubit state \( |\psi\rangle \), the expectation values are defined as
\[
\langle \sigma_i \otimes \sigma_i \rangle
=
\langle \psi|
(\sigma_i \otimes \sigma_i)
|\psi\rangle,
\qquad i\in\{x,y,z\}.
\]

\paragraph{Analytical behavior under the reduced phase model.}

Within the reduced phase model, the initial Bell state
\[
|\Psi^+\rangle = \frac{|01\rangle + |10\rangle}{\sqrt{2}}
\]
evolves into
\[
|\psi(\theta)\rangle
=
\frac{|01\rangle + e^{i\theta}|10\rangle}{\sqrt{2}}.
\]

For this family of states, the correlation observables can be computed analytically:
\[
\langle \sigma_x \otimes \sigma_x \rangle = \cos\theta,
\qquad
\langle \sigma_y \otimes \sigma_y \rangle = \cos\theta,
\qquad
\langle \sigma_z \otimes \sigma_z \rangle = -1.
\]

These results show that the fiber-induced phase \(\theta\) affects only the transverse correlations in the \(x\)- and \(y\)-directions, while the longitudinal correlation along \(z\) remains invariant. This reflects the fact that the relative phase modifies coherence between computational-basis components without changing their populations.

\paragraph{Simulation notes.}
At implementation level, the simulator contains routines for:
\begin{itemize}
    \item generating the Pauli matrices,
    \item constructing tensor-product operators,
    \item evaluating expectation values,
    \item packaging correlation outputs into structured data containers.
\end{itemize}

Additional helper routines verify normalization of input state vectors and unitarity of operators used in the fiber model.

These observables constitute the primary diagnostic outputs of the simulator and provide the most direct bridge between analytical predictions and numerical results. Their dependence on \(\theta\) defines the first validation pipeline of the code.

\subsection{From Deterministic Evolution to Ensemble Description}

The reduced phase model introduced in the previous sections describes a deterministic unitary evolution of a pure entangled state, parametrized by a single phase \(\theta\).

However, in realistic physical conditions, this phase is not constant but fluctuates due to environmental effects such as birefringence variations and thermal instability in the fiber.

If these fluctuations occur on timescales shorter than the measurement resolution, the system is no longer described by a single pure state, but by an ensemble of states corresponding to different phase realizations.

In this regime, the appropriate mathematical description is given by a density operator obtained by averaging over the ensemble:
\[
\rho = \mathbb{E}_{\theta}\left[ |\psi(\theta)\rangle\langle \psi(\theta)| \right].
\]

This transition marks the passage from a deterministic unitary description at the level of each phase realization to an effective mixed-state description obtained after averaging over unresolved phase fluctuations. It provides the conceptual foundation for the second numerical experiment.
For mixed states described by a density operator \(\rho\), the expectation value generalizes to
\[
\langle O \rangle = \mathrm{Tr}(\rho O).
\]

This formulation is essential when dealing with ensemble descriptions, where the system cannot be represented by a single state vector.

%========================================================
%========================================================
\section{Quantitative Characterization of Bipartite States}

\subsection{Reduced Density Operators}

In bipartite quantum systems, the state of an individual subsystem is generally described by a density operator obtained through a partial trace over the complementary subsystem. This formalism accounts for correlations that may be present at the global level.

Let the composite system be defined on the Hilbert space
\[
\mathcal{H}_{AB} = \mathcal{H}_A \otimes \mathcal{H}_B,
\]
with $\mathcal{H}_A \cong \mathbb{C}^2$ and $\mathcal{H}_B \cong \mathbb{C}^2$. A generic global state is described by a density operator
\[
\rho_{AB} \in \mathbb{C}^{4 \times 4}.
\]

The reduced density operators are defined as
\[
\rho_A = \mathrm{Tr}_B(\rho_{AB}), \qquad
\rho_B = \mathrm{Tr}_A(\rho_{AB}).
\]

These operators reproduce all expectation values of observables acting locally on the subsystem. For any observable $O_A$ acting on subsystem $A$,
\[
\langle O_A \rangle = \mathrm{Tr}(\rho_{AB}(O_A \otimes I_B)) = \mathrm{Tr}(\rho_A O_A).
\]

An explicit construction can be written as
\[
\rho_A = \sum_{j=0}^{1} (I_A \otimes \langle j|)\,\rho_{AB}\,(I_A \otimes |j\rangle),
\]
which shows that the partial trace eliminates the degrees of freedom associated with the unobserved subsystem.

For maximally entangled Bell states, such as
\[
|\Psi^+\rangle = \frac{|01\rangle + |10\rangle}{\sqrt{2}},
\]
one obtains
\[
\rho_A = \rho_B = \frac{I}{2},
\]
indicating that each subsystem is locally maximally mixed despite the global state being pure.

The same result holds for the phase-evolved state
\[
|\psi(\theta)\rangle = \frac{|01\rangle + e^{i\theta}|10\rangle}{\sqrt{2}},
\]
where the relative phase modifies global coherence but leaves local states unchanged:
\[
\rho_A(\theta) = \rho_B(\theta) = \frac{I}{2}.
\]

A detailed derivation of this result, including the explicit matrix representation and term-by-term partial trace, is provided in Appendix~A.

By linearity of the partial trace, this property is preserved under convex combinations. Thus, for the ensemble state
\[
\rho_{AB}^{(\mathrm{ens})}
=
\frac{1}{N}
\sum_{k=1}^{N}
|\psi(\theta_k)\rangle\langle\psi(\theta_k)|,
\]
one finds
\[
\rho_A^{(\mathrm{ens})} = \rho_B^{(\mathrm{ens})} = \frac{I}{2}.
\]

This leads to an important conceptual point: the same reduced density operator
\[
\rho_A = \frac{I}{2}
\]
can arise from physically distinct global states.

On one hand, it may describe a classical ensemble
\[
\rho^{(\mathrm{class})}
=
\frac{1}{2}|0\rangle\langle 0|
+
\frac{1}{2}|1\rangle\langle 1|,
\]
where the system is in a definite state but the observer lacks knowledge.

On the other hand, it may result from entanglement:
\[
\rho_A^{(\mathrm{ent})}
=
\mathrm{Tr}_B\big(|\Psi^+\rangle\langle\Psi^+|\big)
=
\frac{I}{2}.
\]

Since both cases yield identical expectation values for all local observables,
\[
\mathrm{Tr}(\rho^{(\mathrm{class})} O_A)
=
\mathrm{Tr}(\rho_A^{(\mathrm{ent})} O_A),
\]
they are indistinguishable by local measurements alone.

The distinction emerges only at the global level, where coherence and correlations must be analyzed through quantities such as purity and bipartite correlation functions (see Appendix~A for an explicit comparison).

Finally, the reduced density operator directly determines measurable local quantities. For Pauli observables,
\[
\langle \sigma_x \rangle = \mathrm{Tr}(\rho_A \sigma_x), \quad
\langle \sigma_y \rangle = \mathrm{Tr}(\rho_A \sigma_y), \quad
\langle \sigma_z \rangle = \mathrm{Tr}(\rho_A \sigma_z),
\]
and for $\rho_A = I/2$ all expectation values vanish, indicating complete local depolarization.

\subsection{Purity and Mixedness of Quantum States}

Having introduced reduced density operators in Section 4.1, we now require a quantitative criterion to distinguish whether a quantum state is pure or mixed. This distinction is essential for interpreting the physical effects of fiber propagation, especially when extending the model beyond deterministic unitary evolution.

A quantum state described by a density operator $\rho$ is said to be \emph{pure} if it can be written as
\[
\rho = |\psi\rangle \langle \psi|,
\]
for some normalized state vector $|\psi\rangle$. Otherwise, the state is said to be \emph{mixed}, meaning that it cannot be represented by a single state vector but only as a statistical ensemble
\[
\rho = \sum_k p_k |\psi_k\rangle \langle \psi_k|, \qquad p_k \ge 0, \quad \sum_k p_k = 1.
\]

A convenient and operational way to distinguish between these two cases is provided by the \emph{purity}, defined as
\[
\gamma(\rho) = \operatorname{Tr}(\rho^2).
\]

This quantity satisfies the fundamental properties
\[
\gamma(\rho) = 1 \quad \Longleftrightarrow \quad \rho \text{ is pure},
\]
\[
\gamma(\rho) < 1 \quad \Longleftrightarrow \quad \rho \text{ is mixed}.
\]

This follows directly from the spectral decomposition of $\rho$. If
\[
\rho = \sum_i \lambda_i |i\rangle \langle i|,
\]
with $\lambda_i \ge 0$ and $\sum_i \lambda_i = 1$, then
\[
\operatorname{Tr}(\rho^2) = \sum_i \lambda_i^2.
\]
The purity reaches its maximum value $1$ only when a single eigenvalue is equal to $1$ and all others vanish, corresponding to a pure state. Otherwise, the sum of squares is strictly smaller than $1$.

More generally, for a system with Hilbert space dimension $d$, the purity is bounded by
\[
\frac{1}{d} \le \operatorname{Tr}(\rho^2) \le 1,
\]
where the lower bound is attained by the maximally mixed state $\rho = I/d$.

\medskip

For single-qubit systems, purity admits a particularly transparent interpretation in terms of the Bloch vector representation introduced in Section 4.1. Any qubit density operator can be written as
\[
\rho = \frac{1}{2}\left(I + \vec{r} \cdot \vec{\sigma}\right),
\]
where $\vec{r} \in \mathbb{R}^3$ is the Bloch vector and $\vec{\sigma} = (\sigma_x, \sigma_y, \sigma_z)$ is the vector of Pauli matrices. In this representation, one finds
\[
\operatorname{Tr}(\rho^2) = \frac{1 + \|\vec{r}\|^2}{2}.
\]
Therefore,
\[
\|\vec{r}\| = 1 \;\Rightarrow\; \rho \text{ is pure},
\qquad
\|\vec{r}\| = 0 \;\Rightarrow\; \rho = \frac{I}{2} \text{ is maximally mixed}.
\]

This result provides a direct connection between the Bloch/Stokes parameters computed in the simulator and the degree of mixedness of the state.

\medskip

As already emphasized in Section 4.1, mixedness does not necessarily imply classical noise. A state may be mixed either because it belongs to a classical statistical ensemble or because it arises as the reduced state of a larger entangled system. In the present context, the latter mechanism is especially important.

A paradigmatic example is provided by the Bell state
\[
|\Psi^+\rangle = \frac{|01\rangle + |10\rangle}{\sqrt{2}}.
\]
Its global density operator
\[
\rho_{AB} = |\Psi^+\rangle \langle \Psi^+|
\]
is pure, and therefore satisfies
\[
\operatorname{Tr}(\rho_{AB}^2) = 1.
\]
However, its reduced density operators are
\[
\rho_A = \rho_B = \frac{I}{2},
\]
which are maximally mixed and satisfy
\[
\operatorname{Tr}(\rho_A^2) = \operatorname{Tr}(\rho_B^2) = \frac{1}{2}.
\]

Thus, a globally pure entangled state may generate locally mixed subsystem states. This distinction between global purity and local mixedness is central to the interpretation of bipartite entangled systems.

\medskip

In the context of the present model, this distinction plays a crucial role. For the phase-evolved Bell-like state
\[
|\psi(\theta)\rangle = \frac{|01\rangle + e^{i\theta}|10\rangle}{\sqrt{2}},
\]
the evolution is generated by local unitary transformations, and therefore the global state remains pure for all values of $\theta$:
\[
\operatorname{Tr}(\rho(\theta)^2) = 1.
\]
At the same time, as shown in Section 4.1, the reduced states remain
\[
\rho_A(\theta) = \rho_B(\theta) = \frac{I}{2},
\]
so that their purity is constant and minimal:
\[
\operatorname{Tr}(\rho_A(\theta)^2) = \operatorname{Tr}(\rho_B(\theta)^2) = \frac{1}{2}.
\]

Therefore, the introduction of a relative phase modifies the correlation observables without affecting either global purity or local mixedness.

\medskip

The notion of purity becomes particularly relevant when extending the model to include more realistic physical effects. 
While deterministic unitary evolution preserves purity, ensemble averaging over random parameters or the inclusion of frequency-dependent propagation effects, such as polarization mode dispersion, can lead to genuine mixed states. 

In particular, for the random-phase ensemble introduced above, the global density operator depends on the ensemble coherence
\[
\mu = \langle e^{i\theta} \rangle,
\]
and its purity is given by
\[
\operatorname{Tr}(\rho_{AB}^{(\mathrm{ens})\,2})
=
\frac{1 + |\mu|^2}{2}.
\]
Equivalently,
\[
\operatorname{Tr}(\rho_{AB}^{(\mathrm{ens})\,2})
=
\frac{1 + \langle \cos\theta \rangle^2 + \langle \sin\theta \rangle^2}{2}.
\]

This shows that the loss of purity is directly governed by the reduction of phase coherence in the ensemble. When $|\mu|=1$, the state remains pure, while for $|\mu|<1$ the ensemble becomes mixed.

For this reason, purity will serve as a fundamental diagnostic tool in the numerical simulator, complementing correlation observables and preparing the ground for more advanced quantities such as fidelity and entanglement measures.

\subsection{Fidelity with Respect to Bell States}

Having introduced reduced density operators and purity as measures of local information and mixedness, we now require a quantitative tool to evaluate how close a given quantum state is to a reference target state. This role is fulfilled by the notion of fidelity.

In general, the fidelity between two quantum states described by density operators $\rho$ and $\sigma$ is defined as
\[
F(\rho,\sigma)
=
\left(\operatorname{Tr}\sqrt{\sqrt{\rho}\,\sigma\,\sqrt{\rho}}\right)^2.
\]

This quantity satisfies
\[
0 \le F(\rho,\sigma) \le 1,
\]
with $F=1$ if and only if $\rho = \sigma$, and $F=0$ for perfectly distinguishable (orthogonal) states. Fidelity therefore provides a measure of similarity between quantum states.

\medskip

In the present context, the reference state is a pure Bell state, and the fidelity simplifies considerably. Let
\[
\sigma = |\phi\rangle\langle\phi|,
\]
with $|\phi\rangle$ a normalized state vector. Then one obtains (see Appendix~B for a detailed derivation)
\[
F(\rho, |\phi\rangle)
=
\langle \phi | \rho | \phi \rangle.
\]

This expression admits a direct physical interpretation: it corresponds to the probability that the state $\rho$ is found in the state $|\phi\rangle$ upon measurement.

\medskip

In this work, we focus on the Bell state
\[
|\Psi^+\rangle = \frac{|01\rangle + |10\rangle}{\sqrt{2}},
\]
and define the fidelity of a generic state $\rho$ with respect to this target as
\[
F_{\Psi^+}(\rho)
=
\langle \Psi^+ | \rho | \Psi^+ \rangle.
\]

\medskip

We now analyze this quantity for the phase-evolved Bell-like state introduced in Section 3.4:
\[
|\psi(\theta)\rangle = \frac{|01\rangle + e^{i\theta}|10\rangle}{\sqrt{2}},
\]
with associated density operator
\[
\rho(\theta) = |\psi(\theta)\rangle\langle\psi(\theta)|.
\]

The fidelity with respect to $|\Psi^+\rangle$ is given by
\[
F_{\Psi^+}(\rho(\theta))
=
|\langle \Psi^+ | \psi(\theta) \rangle|^2.
\]

A direct computation yields
\[
\langle \Psi^+ | \psi(\theta) \rangle
=
\frac{1}{2}(1 + e^{i\theta}),
\]
and therefore
\[
F_{\Psi^+}(\rho(\theta))
=
\frac{1 + \cos\theta}{2}.
\]

\medskip

This result shows that the relative phase $\theta$ modifies the overlap with the reference Bell state, even though the global state remains pure for all values of $\theta$. In particular,
\[
F_{\Psi^+} = 1 \quad \text{for } \theta = 0,
\qquad
F_{\Psi^+} = 0 \quad \text{for } \theta = \pi.
\]

Thus, fidelity captures variations in global coherence that are not reflected in the purity or in the reduced density operators.

\medskip
For ensemble states arising from random phase fluctuations, described by
\[
\rho_{AB}^{(\mathrm{ens})}
=
\frac{1}{N}\sum_{k=1}^{N}
|\psi(\theta_k)\rangle\langle\psi(\theta_k)|,
\]
the fidelity can be computed using the linearity of the expectation value:
\[
F_{\Psi^+}(\rho_{AB}^{(\mathrm{ens})})
=
\frac{1}{N}
\sum_{k=1}^{N}
F_{\Psi^+}(\rho(\theta_k)).
\]

Introducing the ensemble coherence parameter
\[
\mu = \langle e^{i\theta} \rangle,
\]
one finds
\[
F_{\Psi^+}
=
\frac{1}{2}\left(1 + \mathrm{Re}(\mu)\right)
=
\frac{1}{2}\left(1 + \langle \cos\theta \rangle\right).
\]

\medskip

This shows that fidelity directly quantifies the loss of coherence induced by phase fluctuations. In particular, for a uniform distribution over $[0,2\pi]$, one has $\langle \cos\theta \rangle = 0$, leading to
\[
F_{\Psi^+} = \frac{1}{2}.
\]

\medskip

In summary, fidelity provides a quantitative measure of how closely a given state approximates the target Bell state. Unlike purity, which characterizes mixedness, fidelity is sensitive to coherent deformations of the state and therefore plays a central role in the validation and interpretation of the numerical simulations.

\paragraph{Simulation notes.}
Within the simulator architecture, fidelity is introduced as a general-purpose quantitative tool for comparing quantum states in a basis-independent manner. Given a numerically generated state—either in pure vector form or as a density operator—fidelity provides a scalar measure of proximity to a chosen reference state, typically a Bell state in the present framework. 

From an implementation standpoint, this requires supporting both pure-state overlaps and density-matrix evaluations, allowing the same routine to be applied uniformly across different modeling regimes (deterministic evolution, ensemble averaging, or noisy channels). 

Unlike correlation observables, which probe specific measurement directions, fidelity captures a global property of the state and therefore serves as a compact diagnostic for validating transformations, monitoring deviations from ideal entanglement, and benchmarking different physical models within the simulator pipeline.

%========================================================
%========================================================
\section{Numerical Experiments and Results}

\subsection{Experiment 1: Deterministic Phase Sweep}

\subsubsection{Objective}

The first experiment investigates the effect of a controlled relative phase \(\theta\) on a maximally entangled two-qubit state. The objective is to characterize how both correlation observables and state-based quantities evolve under the reduced unitary fiber model.

\subsubsection{Model}

The initial Bell state is chosen as
\[
|\Psi^{+}\rangle
=
\frac{1}{\sqrt{2}}
\left(
|01\rangle + |10\rangle
\right).
\]

After propagation through the effective phase model, the state becomes
\[
|\psi(\theta)\rangle
=
\frac{1}{\sqrt{2}}
\left(
|01\rangle + e^{i\theta}|10\rangle
\right).
\]

The correlation observables evaluated in the simulator are
\[
c_{xx} = \langle \sigma_x \otimes \sigma_x \rangle,\qquad
c_{yy} = \langle \sigma_y \otimes \sigma_y \rangle,\qquad
c_{zz} = \langle \sigma_z \otimes \sigma_z \rangle.
\]

\subsubsection{Analytical Predictions}

For the state \( |\psi(\theta)\rangle \), direct evaluation gives
\[
c_{xx} = \cos\theta,\qquad
c_{yy} = \cos\theta,\qquad
c_{zz} = -1.
\]

Since the evolution is generated by local unitary transformations, the global state remains pure:
\[
\operatorname{Tr}(\rho(\theta)^2)=1.
\]

The reduced density operators remain maximally mixed:
\[
\rho_A(\theta)=\rho_B(\theta)=\frac{I}{2}, \qquad
\operatorname{Tr}(\rho_A^2)=\operatorname{Tr}(\rho_B^2)=\frac{1}{2}.
\]

The fidelity with respect to the reference Bell state \( |\Psi^+\rangle \) is
\[
F_{\Psi^+}(\rho(\theta))
=
\frac{1+\cos\theta}{2}.
\]

\subsubsection{Numerical Results}

The numerical experiment performs a sweep over \(\theta \in [0,2\pi]\), computing both correlation observables and state-characterization quantities.

\begin{figure}[h]
    \centering
    \includegraphics[width=0.75\linewidth]
    {Images/Experiment_1/correlation_observables_theta.png}
    \caption{
    Numerical evaluation of the correlation observables \(c_{xx}\), \(c_{yy}\), and \(c_{zz}\) as functions of the relative phase \(\theta\).
    }
    \label{fig:experiment1_correlations}
\end{figure}

\begin{figure}[h]
    \centering
    \includegraphics[width=0.75\linewidth]{Images/Experiment_1/experiment_1_theory_vs_experiment.png}
    \caption{
    Comparison between numerical simulation and analytical predictions for the correlation observables.
    }
    \label{fig:experiment1_validation}
\end{figure}

\begin{figure}[h]
    \centering
    \includegraphics[width=0.75\linewidth]{Images/Experiment_1/experiment_1_state_metrics.png}
    \caption{
    Validation of global purity, reduced purities, and fidelity with respect to \( |\Psi^+\rangle \).
    }
    \label{fig:experiment1_state_metrics}
\end{figure}

\subsubsection{Discussion}

The numerical results are in full agreement with the analytical predictions:
\begin{itemize}
    \item \(c_{xx}\) and \(c_{yy}\) follow the expected cosine dependence,
    \item \(c_{zz}\) remains constant at \(-1\),
    \item the global purity remains equal to \(1\),
    \item the reduced purities remain equal to \(1/2\),
    \item the fidelity follows \(F_{\Psi^+}=(1+\cos\theta)/2\).
\end{itemize}

These results confirm that the fiber acts as an effective relative phase shifter. While the global state remains pure and maximally entangled, both the observable correlations and the fidelity with respect to a fixed Bell reference vary continuously with the phase. This demonstrates that purity alone is insufficient to characterize the physical transformation induced by the propagation.
%========================================================
%========================================================

\subsection{Experiment 2: Random Phase Ensemble}

\subsubsection{Objective}

The second experiment extends the deterministic phase model by introducing random phase realizations \(\theta_k\), modeling stochastic birefringence effects in realistic optical fibers.

Instead of a single pure state, the system is described by an ensemble
\[
\{ |\psi(\theta_k)\rangle \}_{k=1}^{N},
\]
and by the corresponding averaged density operator.

The objective is to characterize how ensemble averaging modifies global coherence, correlation observables, purity, and fidelity.

\subsubsection{Ensemble Model}

Each realization is given by
\[
|\psi(\theta_k)\rangle
=
\frac{1}{\sqrt{2}}
\left(
|01\rangle + e^{i\theta_k}|10\rangle
\right),
\]
and the ensemble state is
\[
\rho_{\mathrm{ensemble}}
=
\frac{1}{N}
\sum_{k=1}^{N}
|\psi(\theta_k)\rangle\langle \psi(\theta_k)|.
\]

\subsubsection{Analytical Predictions}

\subsubsection{Analytical Predictions}

For the ensemble state, it is convenient to introduce the ensemble coherence parameter
\[
\mu = \langle e^{i\theta} \rangle
=
\langle \cos\theta \rangle + i \langle \sin\theta \rangle.
\]

The ensemble-averaged correlations satisfy
\[
c_{xx} = \mathrm{Re}(\mu) = \langle \cos\theta \rangle,\qquad
c_{yy} = \mathrm{Re}(\mu) = \langle \cos\theta \rangle,\qquad
c_{zz} = -1.
\]

The global purity is
\[
\operatorname{Tr}(\rho_{\mathrm{ensemble}}^2)
=
\frac{1 + |\mu|^2}{2}
=
\frac{1 + \langle \cos\theta \rangle^2 + \langle \sin\theta \rangle^2}{2}.
\]

The reduced states remain maximally mixed:
\[
\rho_A = \rho_B = \frac{I}{2},
\]
and therefore
\[
\operatorname{Tr}(\rho_A^2)=\operatorname{Tr}(\rho_B^2)=\frac{1}{2}.
\]

The fidelity with respect to \( |\Psi^+\rangle \) is
\[
F_{\Psi^+}
=
\frac{1 + \mathrm{Re}(\mu)}{2}
=
\frac{1 + \langle \cos\theta \rangle}{2}.
\]

\subsubsection{Numerical Results}

Three representative phase distributions are considered:
\begin{itemize}
    \item constant phase,
    \item uniform distribution,
    \item Gaussian distribution.
\end{itemize}

Each case is summarized in a single figure combining:
\begin{itemize}
    \item sampled phase statistics,
    \item pure-state correlation trend,
    \item ensemble correlations,
    \item purity and fidelity metrics.
\end{itemize}

\begin{figure}[h]
    \centering
    \includegraphics[width=0.9\linewidth]{Images/Experiment_2/experiment_2_constant_distribution.png}
    \caption{
    Summary of Experiment 2 for a constant phase distribution. The ensemble reduces to a single deterministic realization.
    }
    \label{fig:experiment2_constant}
\end{figure}

\begin{figure}[h]
    \centering
    \includegraphics[width=0.9\linewidth]{Images/Experiment_2/experiment_2_uniform_distribution.png}
    \caption{
    Summary of Experiment 2 for a uniform phase distribution. Phase averaging suppresses transverse correlations and reduces fidelity.
    }
    \label{fig:experiment2_uniform}
\end{figure}

\begin{figure}[h]
    \centering
    \includegraphics[width=0.9\linewidth]{Images/Experiment_2/experiment_2_gaussian_distribution.png}
    \caption{
    Summary of Experiment 2 for a Gaussian phase distribution. Residual coherence depends on the phase variance.
    }
    \label{fig:experiment2_gaussian}
\end{figure}

\subsubsection{Discussion}

The results show a transition from coherent deterministic behavior to mixed-state ensemble behavior:
\begin{itemize}
    \item for constant \(\theta\), the system reproduces Experiment 1,
    \item for distributed \(\theta\), transverse correlations are reduced,
    \item for broad phase distributions, \( |\langle e^{i\theta}\rangle| \to 0 \), leading to vanishing ensemble coherence.
\end{itemize}

At the same time:
\begin{itemize}
    \item the global purity decreases from \(1\),
    \item the reduced purities remain fixed at \(1/2\),
    \item the fidelity with respect to \( |\Psi^+\rangle \) decreases.
\end{itemize}

\subsubsection{Physical Interpretation}

Each realization remains a pure maximally entangled state. The mixedness of the ensemble arises exclusively from statistical averaging over unresolved phase fluctuations.

This demonstrates the transition from a pure-state description to an effective mixed-state description, where coherence is progressively lost at the observable level.

Importantly, the reduction of correlations and fidelity does not necessarily imply complete loss of entanglement, but rather a reduction of ensemble phase coherence accessible through the chosen observables.

%========================================================
%========================================================
\appendix
\section{Explicit Derivation of Reduced Density Operators and Comparison of Global States}

This appendix provides an explicit derivation of the reduced density operators for the phase-evolved Bell-like state and illustrates, through a concrete comparison, why identical local reduced states may correspond to different global physical situations.

We consider the two-qubit state
\[
|\psi(\theta)\rangle = \frac{|01\rangle + e^{i\theta}|10\rangle}{\sqrt{2}},
\]
written in the computational basis
\[
\{|00\rangle, |01\rangle, |10\rangle, |11\rangle\}.
\]

Its vector representation is
\[
|\psi(\theta)\rangle =
\frac{1}{\sqrt{2}}
\begin{pmatrix}
0 \\
1 \\
e^{i\theta} \\
0
\end{pmatrix}.
\]

The associated density operator is
\[
\rho_{AB}(\theta) = |\psi(\theta)\rangle\langle\psi(\theta)|,
\]
which expands as
\[
\rho_{AB}(\theta)
=
\frac{1}{2}
\Big(
|01\rangle\langle 01|
+
e^{-i\theta}|01\rangle\langle 10|
+
e^{i\theta}|10\rangle\langle 01|
+
|10\rangle\langle 10|
\Big).
\]

In matrix form,
\[
\rho_{AB}(\theta)
=
\frac{1}{2}
\begin{pmatrix}
0 & 0 & 0 & 0 \\
0 & 1 & e^{-i\theta} & 0 \\
0 & e^{i\theta} & 1 & 0 \\
0 & 0 & 0 & 0
\end{pmatrix}.
\]

We now compute the reduced density operator of subsystem $A$ by taking the partial trace over subsystem $B$:
\[
\rho_A(\theta) = \mathrm{Tr}_B(\rho_{AB}(\theta)).
\]

Using the explicit formula
\[
\rho_A(\theta)
=
\sum_{b=0}^{1}
(I_A \otimes \langle b|)\,\rho_{AB}(\theta)\,(I_A \otimes |b\rangle),
\]
we evaluate the contribution of each term in $\rho_{AB}(\theta)$ separately.

For the diagonal term $|01\rangle\langle 01|$, with
\[
|01\rangle = |0\rangle_A \otimes |1\rangle_B,
\]
one obtains
\[
\mathrm{Tr}_B(|01\rangle\langle 01|) = |0\rangle\langle 0|.
\]

For the coherence term
\[
|01\rangle\langle 10|
=
(|0\rangle\langle 1|)\otimes(|1\rangle\langle 0|),
\]
one finds
\[
\mathrm{Tr}_B(|01\rangle\langle 10|)
=
|0\rangle\langle 1|\,\mathrm{Tr}(|1\rangle\langle 0|)
=
0,
\]
since
\[
\mathrm{Tr}(|1\rangle\langle 0|)=\langle 0|1\rangle = 0.
\]
The same argument gives
\[
\mathrm{Tr}_B(|10\rangle\langle 01|) = 0.
\]

Finally, for the second diagonal term $|10\rangle\langle 10|$, with
\[
|10\rangle = |1\rangle_A \otimes |0\rangle_B,
\]
one obtains
\[
\mathrm{Tr}_B(|10\rangle\langle 10|) = |1\rangle\langle 1|.
\]

Summing all contributions, the reduced density operator of subsystem $A$ is
\[
\rho_A(\theta)
=
\frac{1}{2}|0\rangle\langle 0|
+
\frac{1}{2}|1\rangle\langle 1|
=
\frac{I}{2}.
\]

By symmetry, the same result holds for subsystem $B$:
\[
\rho_B(\theta) = \frac{I}{2}.
\]

Therefore, although the global density operator depends explicitly on the phase $\theta$, the reduced density operators do not. The phase appears only in off-diagonal coherence terms, and these are eliminated by the partial trace because of the orthogonality relations in the traced subsystem.

\medskip

Setting $\theta = 0$ recovers the Bell state
\[
|\Psi^+\rangle = \frac{|01\rangle + |10\rangle}{\sqrt{2}},
\]
which gives the same reduced density operators
\[
\rho_A = \rho_B = \frac{I}{2}.
\]

This shows explicitly that each subsystem is locally maximally mixed even though the global state is pure.

\medskip

To understand why this does not imply classical ignorance, consider the single-qubit mixed state
\[
\rho^{(\mathrm{class})}
=
\frac{1}{2}|0\rangle\langle 0|
+
\frac{1}{2}|1\rangle\langle 1|
=
\frac{I}{2}.
\]

At the level of subsystem $A$, the reduced state obtained from the entangled two-qubit state coincides with this classical mixture:
\[
\rho_A^{(\mathrm{ent})} = \rho^{(\mathrm{class})}.
\]

Consequently, for any local observable $O_A$,
\[
\mathrm{Tr}(\rho^{(\mathrm{class})} O_A)
=
\mathrm{Tr}(\rho_A^{(\mathrm{ent})} O_A).
\]

Thus, no measurement acting only on subsystem $A$ can distinguish whether the local mixedness arises from a classical mixture or from entanglement with subsystem $B$.

\medskip

At the global level, however, the distinction becomes manifest. Consider the separable mixed state
\[
\rho_{AB}^{(\mathrm{sep})}
=
\frac{1}{2}|01\rangle\langle 01|
+
\frac{1}{2}|10\rangle\langle 10|.
\]

This state has the same reduced density operators as the entangled state,
\[
\mathrm{Tr}_B(\rho_{AB}^{(\mathrm{sep})}) = \frac{I}{2},
\qquad
\mathrm{Tr}_A(\rho_{AB}^{(\mathrm{sep})}) = \frac{I}{2},
\]
but it does not contain the off-diagonal coherence terms
\[
e^{-i\theta}|01\rangle\langle 10|,
\qquad
e^{i\theta}|10\rangle\langle 01|.
\]

One immediate indicator of this difference is the purity
\[
\mathrm{Tr}(\rho^2).
\]

For the entangled pure state,
\[
\mathrm{Tr}(\rho_{AB}(\theta)^2) = 1,
\]
while for the separable mixed state,
\[
\mathrm{Tr}\big((\rho_{AB}^{(\mathrm{sep})})^2\big) = \frac{1}{2}.
\]

A second distinction appears in the bipartite correlation observables. For the phase-evolved state,
\[
\langle \sigma_x \otimes \sigma_x \rangle = \cos\theta, \qquad
\langle \sigma_y \otimes \sigma_y \rangle = \cos\theta, \qquad
\langle \sigma_z \otimes \sigma_z \rangle = -1,
\]
whereas for the separable mixed state one obtains
\[
\langle \sigma_x \otimes \sigma_x \rangle = 0, \qquad
\langle \sigma_y \otimes \sigma_y \rangle = 0, \qquad
\langle \sigma_z \otimes \sigma_z \rangle = -1.
\]

Thus, the two states are identical from the local point of view but differ at the global level in both purity and correlation structure. This confirms explicitly that reduced density operators encode all locally accessible information, but do not in general capture global quantum coherence or bipartite entanglement.

%==============================================================
\section{Derivation of Fidelity for a Pure Reference State}

We derive the simplified expression for the fidelity when one of the states is pure, and justify the algebraic properties used in the reduction.

\medskip

Let $\rho$ be a density operator on a finite-dimensional Hilbert space, and let
\[
\sigma = |\psi\rangle\langle\psi|
\]
be a pure-state density operator, with $|\psi\rangle$ normalized:
\[
\langle \psi | \psi \rangle = 1.
\]

By construction, $\sigma$ is Hermitian,
\[
\sigma^\dagger = \sigma,
\]
positive semidefinite,
\[
\langle \chi | \sigma | \chi \rangle \ge 0
\quad \text{for all } |\chi\rangle,
\]
and satisfies
\[
\operatorname{Tr}(\sigma) = 1.
\]
Moreover, it is a rank-one operator, since it projects onto the one-dimensional subspace spanned by $|\psi\rangle$.

\medskip

The general definition of fidelity is
\[
F(\rho,\sigma)
=
\left(\operatorname{Tr}\sqrt{\sqrt{\rho}\,\sigma\,\sqrt{\rho}}\right)^2.
\]

Substituting $\sigma = |\psi\rangle\langle\psi|$, one obtains
\[
\sqrt{\rho}\,\sigma\,\sqrt{\rho}
=
\sqrt{\rho}\,|\psi\rangle\langle\psi|\,\sqrt{\rho}.
\]

Define the vector
\[
|\phi\rangle = \sqrt{\rho}\,|\psi\rangle.
\]
Then the operator becomes
\[
\sqrt{\rho}\,\sigma\,\sqrt{\rho}
=
|\phi\rangle\langle\phi|.
\]

This is again a rank-one operator, provided $|\phi\rangle \neq 0$.

\medskip

We now compute the square root of this operator. Since
\[
|\phi\rangle\langle\phi|
\]
is positive semidefinite and rank one, its only nonzero eigenvalue is
\[
\lambda = \langle \phi | \phi \rangle.
\]

Thus, its spectral decomposition is
\[
|\phi\rangle\langle\phi|
=
\lambda\,|\hat{\phi}\rangle\langle\hat{\phi}|,
\]
where
\[
|\hat{\phi}\rangle = \frac{|\phi\rangle}{\|\phi\|}.
\]

The square root is then
\[
\sqrt{|\phi\rangle\langle\phi|}
=
\sqrt{\lambda}\,|\hat{\phi}\rangle\langle\hat{\phi}|.
\]

Taking the trace,
\[
\operatorname{Tr}\sqrt{|\phi\rangle\langle\phi|}
=
\sqrt{\lambda}.
\]

Therefore,
\[
F(\rho,\sigma)
=
\lambda
=
\langle \phi | \phi \rangle.
\]

Recalling the definition of $|\phi\rangle$,
\[
\lambda
=
\langle \psi | \rho | \psi \rangle.
\]

\medskip

We conclude that when $\sigma = |\psi\rangle\langle\psi|$ is a pure state, the fidelity reduces to
\[
F(\rho,|\psi\rangle\langle\psi|)
=
\langle \psi | \rho | \psi \rangle.
\]

\medskip

This result justifies the simplified expression used in Section 4.3. It also shows that the fidelity corresponds to the expectation value of the projector onto $|\psi\rangle$, providing a direct probabilistic interpretation as the probability of obtaining the outcome $|\psi\rangle$ in a projective measurement.
%========================================================
%========================================================
%========================================================
\section{References}

\begin{itemize}
    \item M. A. Nielsen and I. L. Chuang, \textit{Quantum Computation and Quantum Information}, Cambridge University Press.
\end{itemize}

\end{document}